\documentclass[12pt,a4paper]{article}
\usepackage[utf8]{inputenc}
\usepackage[spanish,es-nodecimaldot]{babel}
\usepackage{amsmath}
\usepackage{amssymb}
\usepackage{graphicx}
\usepackage{booktabs}
\usepackage[hidelinks]{hyperref}
\usepackage{natbib}
\usepackage{geometry}
\usepackage{setspace}
\usepackage{caption}
\usepackage{subcaption}
\usepackage{array}
\geometry{margin=2.5cm}
\onehalfspacing

\title{\textbf{Efecto de variables morfométricas sobre la ocurrencia
de movimientos en masa en el Valle de Aburrá mediante\\
Modelos Aditivos Generalizados (GAM)}}
\author{Edier Aristizábal}

\begin{document}

\maketitle

%-------------------------------------------------------------------
\begin{abstract}
%-------------------------------------------------------------------
Los modelos estadísticos para la evaluación de la susceptibilidad por
movimientos en masa frecuentemente asumen relaciones lineales entre las
variables morfométricas y la probabilidad de ocurrencia, o bien
emplean modelos no paramétricos de difícil interpretación. Este
estudio aplica Modelos Aditivos Generalizados (GAM) para cuantificar
el efecto no lineal de seis variables morfométricas —elevación,
pendiente, aspecto, curvatura planiforme, curvatura de perfil e
índice topográfico de humedad (TWI)— sobre la ocurrencia de
movimientos en masa en el Valle de Aburrá, Colombia. Se utilizó un
inventario de 5\,487 eventos históricos (2001--2024) y un DEM de
2\,m de resolución. El modelo, evaluado en un conjunto de prueba
independiente (partición 70/30), alcanzó AUC = 0{,}908 y TSS =
0{,}816. El TWI fue el predictor dominante, seguido de la pendiente
y las curvaturas topográficas. Las funciones de suavizado permiten
identificar rangos críticos de cada variable morfométrica asociados
a la mayor probabilidad de falla, constituyendo información de alto
valor para la gestión del riesgo y el ordenamiento territorial en el
área metropolitana.

\bigskip\noindent\textbf{Palabras clave:} movimientos en masa,
susceptibilidad, GAM, \textit{splines}, Valle de Aburrá,
geomorfometría, TWI.
\end{abstract}

%===================================================================
\section{Introducción}
%===================================================================

Los movimientos en masa constituyen una de las amenazas naturales de
mayor impacto humano y económico a escala global, especialmente en
regiones tropicales de montaña donde confluyen precipitaciones intensas
y vertientes pronunciadas \citep{aristizabal2026landslide,
gomez2023spatial, petley2012global}. El análisis de las bases de datos
globales revela que entre 1903 y 2020 se han
documentado más de 37\,946 deslizamientos que produjeron 185\,753
víctimas fatales en 161 países, con los continentes americano y asiático
concentrando las mayores cifras \citep{gomez2023spatial}.
Aproximadamente el 60\,\% de los deslizamientos reportados a nivel
mundial son desencadenados por lluvias, lo que subraya el papel
dominante de la precipitación como mecanismo detonante y la necesidad
de integrar variables hidrometeorológicas y morfométricas en los modelos
de evaluación de la amenaza \citep{gomez2023spatial,
aristizabal2026landslide}.

Colombia se encuentra entre los países más afectados del hemisferio
occidental. El inventario nacional registra 30\,730 deslizamientos
ocurridos entre 1900 y 2018, con más de 34\,000 víctimas fatales
asociadas; el 92\,\% de estos eventos fueron desencadenados por lluvia
y el 80\,\% ocurrieron en los flancos de las cordilleras de lso Andes
\citep{aristizabal2020landslide}. Este patrón espacial refleja la
combinación de laderas de fuerte pendiente, perfiles de meteorización profundos, y una hidroclimatología dominada por el desplazamiento de la
Zona de Convergencia Intertropical (ZCIT), que impone un régimen de
lluvias bimodal con máximas en abril--junio y septiembre--noviembre
\citep{aristizabal2026landslide}. La variabilidad interanual asociada
al ENSO (\textit{El Niño-Southern Oscillation}) amplifica este patrón:
las fases La~Niña intensifican las lluvias durante las temporadas
húmedas, elevando la humedad antecedente del suelo y predisponiendo las
laderas a la falla \citep{aristizabal2026landslide}. 

Los primeros
registros sistemáticos de desastres por movimientos en masa en el Valle
de Aburrá datan de 1880 \citep{aristizabal2007inventario}, y la base
de datos regional acumula más de 6\,750 eventos entre 1880 y 2007,
de los cuales los movimientos en masa representan el 35\,\% del total
de emergencias registradas. En el Valle de Aburrá, con
aproximadamente 4{,}1 millones de habitantes, los movimientos en masa
son responsables del 75\,\% de las víctimas por desastres y
constituyen la amenaza geológica de mayor recurrencia
\citep{aristizabal2020landslide, nieto2025coupled}. El crecimiento
urbano extraordinariamente rápido del valle —la población se incrementó
30~veces entre 1905 y 2005, pasando de 103\,300 a más de 3{,}3
millones de habitantes \citep{nieto2025coupled}— ha llevado a la
ocupación progresiva de laderas cada vez más pronunciadas, en particular
mediante asentamientos informales con infraestructura deficiente de
drenaje y contención \citep{nieto2025coupled, aristizabal2026landslide}.
La evidencia empírica indica que los deslizamientos causan entre
2 y 3{,}5~veces más víctimas por km² en barrios informales que en los
sectores formales de mayor estrato \citep{nieto2025coupled}, y que el
reciente incremento en el número e impacto de los eventos está asociado
principalmente a la expansión urbana hacia terrenos intrínsecamente más
inestables, más que a cambios en el régimen de lluvias
\citep{nieto2025coupled, aristizabal2026landslide}.

La evaluación de la susceptibilidad, definida como la condición intrínseca del
terreno a generar movimientos en masa, es el primer paso en la
cadena de análisis de riesgo \citep{hungr2014varnes}. Las metodologías para evaluar la susceptibilidad por movimientos en masa se clasifican en: (1)~\textit{heurísticos}, basados en juicio
experto y ponderación de mapas temáticos; (2)~\textit{físicos}, que
resuelven ecuaciones de estabilidad pero requieren parámetros
geomecánicos de alta resolución difíciles de obtener regionalmente; y
(3)~\textit{estadísticos}, que infieren la susceptibilidad a partir de
la relación histórica entre la distribución espacial de los eventos y
las covariables del terreno \citep{brenning2008statistical}.

Dentro de los métodos estadísticos, la regresión logística
(RL) ha sido el estándar por su interpretabilidad y solidez teórica.
Sin embargo, la RL impone relaciones lineales entre el predictor lineal
y las covariables en escala logit. Por ejemplo, la
susceptibilidad aumenta con la pendiente hasta un rango óptimo y luego
disminuye donde el material susceptible ya ha sido removido por
episodios previos de inestabilidad. En el otro extremo, los métodos de apreddizaje automático (\textit{machine learning} - bosques aleatorios, redes neuronales,
máquinas de gradiente) superan esta restricción, pero producen modelos
de caja negra cuyas predicciones son difícilmente validables desde el
punto de vista geomorfológico y cuya utilidad para identificar los
factores condicionantes físicos es limitada
\citep{merghadi2020machine}.

Los Modelos Aditivos Generalizados (GAM)
\citep{hastie1990generalized, wood2017generalized} ofrecen una solución
que equilibra ambas ventajas: permiten relaciones no lineales arbitrarias entre
covariables y respuesta mediante funciones de suavizado
(\textit{splines}), al tiempo que conservan la interpretabilidad
gráfica y estadística de los modelos lineales generalizados. Cada
función de suavizado puede graficarse individualmente, lo que permite
cuantificar y comunicar el efecto de cada variable morfométrica sobre
la probabilidad de ocurrencia de movimientos en masa.

El presente trabajo tiene como objetivo principal cuantificar
el efecto no lineal de seis variables morfométricas sobre la ocurrencia
de movimientos en masa en el Valle de Aburrá, utilizando GAM con
regresión logística como marco estadístico. Aunque la generación del mapa de
susceptibilidad es un producto derivado del análisis, no corresponde al objetivo del presente estudio, por lo cual no se incorporaron dentro del análisis otras variables importantes en la ocurrencia d movimientos en masa como geología, coberturas del suelo, entre otros.

%===================================================================
\section{Área de Estudio}
%===================================================================

El Valle de Aburrá es un valle intramontano ubicado sobre el tope del norte de la Cordillera Central de los Andes colombianos,
departamento de Antioquia (Figura~\ref{fig:localizacion}). Se
extiende entre los 6°05' y 6°35' de latitud norte y los 75°25' y
75°45' de longitud oeste, con aproximadamente 60\,km de extensión
norte--sur y hasta 10\,km de anchura. Alberga diez municipios del Área
Metropolitana, con una población de aproximadamente
3{,}7 millones de habitantes.

\begin{figure}[htbp]
    \centering
    \includegraphics[width=0.55\textwidth]{FIGURAS/mapa_localizacion.png}
    \caption{Localización del Valle de Aburrá en el contexto de
    Colombia. El recuadro rojo delimita el área de estudio en el
    departamento de Antioquia.}
    \label{fig:localizacion}
\end{figure}

Desde el punto de vista topográfico, el valle es profundamente incisado
por el río Medellín, con cotas que varían entre 1\,300\,m
en el fondo y más de 2\,800\,m en las cimas de las vertientes.
Las pendientes predominantes oscilan entre 20° y 50°, con una
heterogeneidad morfológica marcada por afloramientos en suelos
residuales y depósitos de ladera.

Geológicamente, el valle está enmarcado en el cinturón orogénico
formado por la convergencia oblicua de las placas de Nazca y del
Caribe con la placa Suramericana \citep{trenkamp2002wide}, lo que ha
generado intensa deformación cortical, levantamiento y fallamiento
\citep{cediel2003phanerozoic}. El subsuelo está dominado por rocas
metamórficas (esquistos, gneis y anfibolitas del Complejo Cajamarca)
y cuerpos ígneos intrusivos (granodioritas, cuarzodioritas), cubiertos
por espesos perfiles de meteorización de hasta 30--50\,m. La
predominancia de suelos residuales y depósitos coluviales con
propiedades geomecánicas espacialmente variables —en términos de
textura, cohesión y conductividad hidráulica— modula localmente la
infiltración, el desarrollo de presiones de poros y la estabilidad de
laderas bajo lluvia intensa \citep{aristizabal2026landslide,
nieto2025coupled}.

El clima es subtropical de montaña, controlado por la migración
meridional de la Zona de Convergencia Intertropical (ZCIT) y la
interacción de tres chorros de bajo nivel: el Chocó, el del Caribe y
el de los Andes orientales \citep{aristizabal2026landslide}. El
régimen pluviométrico es bimodal, con temporadas húmedas en
marzo--mayo y septiembre--noviembre, precipitación media anual entre
1\,500 y 3\,000\,mm, y temperatura entre 12 y 24\,°C según la
altitud. La variabilidad interanual está modulada principalmente por
el ENSO: las fases La~Niña intensifican las lluvias de ambas
temporadas húmedas, elevando la humedad antecedente del suelo y
predisponiendo las laderas a la falla
\citep{aristizabal2026landslide}.

%===================================================================
\section{Datos y Metodología}
%===================================================================

\subsection{Marco conceptual: los Modelos Aditivos Generalizados (GAM)}

La susceptibilidad por movimientos en masa depende de variables morfométricas
que raramente exhiben relaciones lineales con la probabilidad de falla.
Un modelo de regresión logística clásica asumiría, por ejemplo, que a mayor
pendiente la probabilidad de deslizamiento aumenta en proporción constante;
sin embargo, la evidencia empírica muestra que pendientes muy bajas
($<\!5^{\circ}$) y muy altas ($>\!55^{\circ}$) presentan baja frecuencia de eventos,
con un rango de mayor ocurrencia situado entre 20° y 45°.
Relaciones en forma de cúpula, con umbrales o con inversiones de efecto
en distintos rangos del predictor no pueden ser representadas por funciones
lineales. Otro ejemplo relevante es el TWI: su efecto sobre la probabilidad
de falla no es monótonamente creciente sino que se satura o invierte a
valores muy altos cuando el suelo permanece permanentemente saturado y los
materiales finos han sido lavados. Modelar incorrectamente estas no
linealidades sesga las predicciones espaciales de susceptibilidad y
puede llevar a errar en la identificación de zonas críticas.

Los Modelos Aditivos Generalizados
\citep{hastie1990generalized, wood2017generalized} extienden los
Modelos Lineales Generalizados (GLM) reemplazando los efectos lineales por
\textit{funciones de suavizado} no paramétricas. Cada predictor contribuye
al modelo mediante una curva flexible, estimada directamente de los datos,
sin imponer ninguna forma funcional predefinida. El modelo \textit{suma}
estas contribuciones individuales ---de allí el calificativo
``aditivo''--- lo que preserva su interpretabilidad: el efecto parcial de
cada variable puede visualizarse y analizarse de forma independiente,
controlando simultáneamente por todas las demás.

Dado que la variable respuesta es binaria (presencia/ausencia de
movimiento en masa) se asume $Y_i \sim \mathrm{Bernoulli}(P_i)$ y se
adopta la función de enlace logit para modelar la probabilidad de
ocurrencia $P_i$ en cada celda del territorio:

\begin{equation}
\ln\!\left(\frac{P_i}{1-P_i}\right) = \beta_0 +
s_1(\mathrm{elev}_i) + s_2(\mathrm{slope}_i) +
s_3(\mathrm{aspect}_i) + s_4(\kappa_{p,i}) +
s_5(\kappa_{r,i}) + s_6(\mathrm{TWI}_i),
\label{eq:gam}
\end{equation}

donde $\beta_0$ es el intercepto global y $s_j(\cdot)$ son funciones
de suavizado que capturan la relación no lineal entre cada predictor y
el logit de la probabilidad (el logaritmo de la razón entre la
probabilidad de ocurrencia y la de no-ocurrencia). La probabilidad se
recupera aplicando la función logística inversa:
$P_i = [1 + \exp(-\eta_i)]^{-1}$, donde $\eta_i$ es el predictor
lineal (parte derecha de la ecuación~\eqref{eq:gam}).

\subsubsection*{Construcción de las funciones de suavizado}

Cada función $s_j$ se representa matemáticamente como combinación
lineal de $K$ funciones de base $\{\phi_k\}_{k=1}^{K}$:
$s_j(x) = \sum_{k=1}^{K} \gamma_{jk}\,\phi_k(x)$.
Las funciones de base actúan como ``ladrillos'' de distintas
formas (curvas elementales de diferentes frecuencias); combinándolas
con distintos pesos $\gamma_{jk}$ es posible aproximar cualquier curva
suave. $K$ es la \textit{dimensión de la base} y controla la
complejidad máxima de la curva.

En este estudio se emplean \textit{thin plate regression splines}
(TPRS) para las cinco variables continuas no circulares
\citep{wood2017generalized}. Los TPRS ubican automáticamente los nodos
en posiciones óptimas (no se requiere especificar su posición a priori)
y, para una restricción de rugosidad dada, minimizan el error cuadrático
de interpolación. Esto los hace teóricamente óptimos entre los
\textit{splines} de dimensión $K$ dada.

\subsubsection*{Penalización de rugosidad y parámetro de suavizado}

Sin restricciones, la curva $s_j$ podría ajustarse perfectamente a los
datos de entrenamiento mediante oscilaciones erráticas de escaso
significado físico (sobreajuste). Para evitarlo, se penaliza la
rugosidad de la función mediante la integral del cuadrado de su segunda
derivada:

\begin{equation}
    J(s_j) = \int \bigl[s_j''(x)\bigr]^2\,dx.
    \label{eq:penalty}
\end{equation}

La segunda derivada $s_j''(x)$ mide la curvatura local de la función:
si $s_j''(x) = 0$ en todo el dominio, la función es una línea recta;
valores altos de $|s_j''|$ indican cambios abruptos de dirección.
La integral $J(s_j)$ acumula esa curvatura a lo largo del dominio
completo del predictor: es nula para una recta y crece con la
ondulación de la curva. La estimación del modelo maximiza la
\textit{log-verosimilitud penalizada}, que balancea la bondad del
ajuste con la suavidad de las funciones:

\begin{equation}
    \ell_P = \ell(\boldsymbol{\gamma}) -
    \tfrac{1}{2}\sum_j \lambda_j\, J(s_j),
    \label{eq:plik}
\end{equation}

donde $\ell(\boldsymbol{\gamma})$ es la log-verosimilitud ordinaria
(cuán bien el modelo reproduce los datos observados) y
$\lambda_j > 0$ es el \textit{parámetro de suavizado} del término $j$.
Este parámetro gobierna el compromiso entre ajuste y suavidad: cuando
$\lambda_j \to 0$ la penalización desaparece y la curva puede ser
muy flexible (riesgo de sobreajuste); cuando $\lambda_j \to \infty$
la penalización domina y la curva converge a una línea recta,
recuperando el caso de la regresión logística clásica. Los GAM
incluyen, por tanto, a los GLM como caso límite particular.

\subsubsection*{Selección automática del suavizado: REML}

La determinación del valor apropiado de $\lambda_j$ es crítica para el
desempeño del modelo: valores demasiado pequeños producen sobreajuste
y pérdida de generalización; valores demasiado grandes generan
infraajuste y eliminan las relaciones no lineales de interés.
En este estudio, los $\lambda_j$ se seleccionaron automáticamente
mediante \textit{Restricted Maximum Likelihood} (REML)
\citep{wood2017generalized}. REML reinterpreta el GAM como un modelo
lineal mixto equivalente, en el que los parámetros de suavizado
corresponden a varianzas de efectos aleatorios. Este enfoque permite
estimar simultáneamente todos los $\lambda_j$ en un único procedimiento
de optimización estadística, y proporciona inferencia rigurosa sobre
los grados de libertad efectivos de cada curva. Frente a la validación
cruzada generalizada (GCV), REML produce estimaciones más estables y
es menos propenso a sobresuavizar cuando los datos presentan correlación
espacial o ligero desequilibrio de clases
\citep{wood2017generalized}, condiciones comunes en inventarios de
movimientos en masa.

\subsection{Datos e inventario de movimientos en masa}

Para el análisis se construyó un inventario de movimientos en masa que comprende 5\,487
puntos de ocurrencia entre 2001 y 2024 (Figura~\ref{fig:inventario}),
a partir de imágenes de satélite de resolución submétrica disponibles en Google Earth.

\begin{figure}[htbp]
    \centering
    \includegraphics[width=0.70\textwidth]{FIGURAS/mapa_inventario.png}
    \caption{Distribución espacial de los 5\,487 movimientos en masa
    inventariados (2001--2024) en el Valle de Aburrá. Puntos rojos
    sobre hillshade derivado del DEM de 5\,m.}
    \label{fig:inventario}
\end{figure}

Se construyó un Modelo Digital de Elevación (MDE) con la mejor resolución
espacial disponible a partir de un modelo de 1\,m de resolución para la
ciudad de Medellín, obtenido con tecnología LiDAR, y un MDE complementario
con resolución de 2\,m para suelos urbanos y 5\,m para suelos rurales del
resto del valle, derivado de restitución fotogramétrica.

A partir de dicho MDE se derivaron seis variables morfométricas (Tabla~\ref{tab:estadisticos}
y Figura~\ref{fig:hist}):

\begin{enumerate}

\item \textbf{Elevación} ($z$, m): altura sobre el nivel del mar.
En el Valle de
Aburrá, los sectores medios de ladera (1\,500--2\,000\,m) presentan
la mayor densidad
constructiva sobre ladera, lo que condiciona una mayor concentración
de eventos en ese rango altitudinal (Figura~\ref{fig:hist}).

\item \textbf{Pendiente} ($\beta$, grados): gradiente del terreno.
Está directamente relacionada con la estabilidad de laderas, dado que
la componente tangencial de la gravedad, que genera tensión de corte
en el plano de falla, aumenta con $\sin\beta$. Las distribuciones de
densidad muestran una concentración marcada de deslizamientos entre 20°
y 45°, con pico alrededor de 30°. La reducción de la frecuencia de
eventos en pendientes superiores a 50° es coherente con la erosión
progresiva del material susceptible y la escasa área relativa de estas
superficies en el valle.

\item \textbf{Aspecto} ($\alpha$, grados): orientación de la ladera
respecto al norte (0°--360°). Determina la exposición solar, los
patrones de evapotranspiración y, como aproximación, la humedad del
suelo. El aspecto es una variable \textit{circular}: un valor de 360°
es físicamente idéntico a 0°. Un TPRS estándar trataría esos dos
extremos como valores opuestos del dominio, generando una
discontinuidad artificial que distorsionaría la función estimada.
Para evitar este artefacto se empleó un \textit{spline} cúbico
cíclico (\texttt{bs\,=\,"cc"}), que impone la continuidad de la
función y su primera derivada en los extremos del dominio circular:
\begin{equation}
    s_3(0^{\circ}) = s_3(360^{\circ}) \quad\text{y}\quad
    s_3'(0^{\circ}) = s_3'(360^{\circ}),
    \label{eq:cyclic}
\end{equation}
con nodos extremos fijados en 0° y 360° y dimensión de base $k=10$.
Esta condición garantiza que la función estimada cierra suavemente
al completar el círculo de orientaciones, sin saltos ni pendientes
discontinuas en la transición norte--norte.

\item \textbf{Curvatura planiforme} ($\kappa_p$): curvatura
perpendicular a la dirección de máxima pendiente. Controla la
convergencia o divergencia lateral del flujo subsuperficial. Valores
negativos (cóncavo en planta) indican vaguadas y hondonadas donde
se concentra la humedad y la presión de poros; valores positivos
corresponden a filos y divisorias donde el agua diverge.

\item \textbf{Curvatura de perfil} ($\kappa_r$): curvatura en la
dirección de máxima pendiente. Controla la aceleración del flujo
ladera abajo y el espesor del material coluvial: zonas cóncavas en
perfil tienden a acumular el flujo de agua, aumentando el peso del
material y la presión de poros, predisponiendo a la falla.

\item \textbf{Índice Topográfico de Humedad (TWI)}: propuesto por
\citet{beven1979physically}, cuantifica la propensión topográfica
a acumular agua:
\begin{equation}
    \mathrm{TWI} = \ln\!\left(\frac{A_s}{\tan\beta}\right),
    \label{eq:twi}
\end{equation}
donde $A_s$ (m²/m) es el área de contribución específica aguas
arriba. Valores altos de TWI indican tendencia a la saturación del
suelo, condición que reduce la resistencia al corte en suelos
cohesivos y facilita la movilización de materiales ladera abajo.

\end{enumerate}

La Tabla~\ref{tab:estadisticos} resume los estadísticos descriptivos
de las seis variables para el conjunto de datos utilizado.

\begin{table}[htbp]
\centering
\caption{Estadísticos descriptivos de las variables morfométricas
del conjunto de datos de entrenamiento ($n = 10\,974$).
DE: desviación estándar.}
\label{tab:estadisticos}
\small
\begin{tabular}{lrrrrrrrr}
\toprule
Variable & Mínimo & Q1 & Media & Mediana & Q3 & Máximo & DE & Moda \\
\midrule
Elevación (m)       & 1\,303 & 1\,551 & 1\,762 & 1\,698 & 1\,962 & 2\,847 & 298 & 1\,632 \\
Pendiente (°)       & 0{,}1  & 19{,}4 & 27{,}3 & 26{,}8 & 34{,}7 & 69{,}4 & 10{,}2 & 28{,}1 \\
Aspecto (°)         & 0{,}0  & 89{,}6 & 181{,}5 & 181{,}2 & 271{,}4 & 360{,}0 & 103{,}4 & 175{,}3 \\
Curv. planiforme    & $-4{,}82$ & $-0{,}31$ & $-0{,}03$ & $-0{,}01$ & 0{,}28 & 4{,}97 & 0{,}87 & $-0{,}02$ \\
Curv. de perfil     & $-4{,}91$ & $-0{,}33$ & $-0{,}02$ & $-0{,}01$ & 0{,}30 & 5{,}02 & 0{,}91 & $-0{,}01$ \\
TWI                 & 1{,}24 & 3{,}62 & 4{,}87 & 4{,}52 & 5{,}97 & 12{,}31 & 1{,}83 & 4{,}18 \\
\bottomrule
\end{tabular}
\end{table}

\begin{figure}[htbp]
    \centering
    \includegraphics[width=0.95\textwidth]{FIGURAS/comparativa_histogramas.png}
    \caption{Distribuciones de densidad de las seis variables
    morfométricas para deslizamientos (rojo, $n = 5\,487$) y
    no-ocurrencias (azul, $n = 5\,487$). Las diferencias son
    estadísticamente significativas en todos los casos
    ($p < 0{,}001$, prueba KS).}
    \label{fig:hist}
\end{figure}

\subsection{Diseño del modelo: muestreo, ajuste y diagnósticos}

\subsubsection*{Muestreo de no-ocurrencias}

Para construir un modelo de presencia--ausencia es necesario definir
puntos de no-ocurrencia representativos del territorio estable.
Se generaron 5\,487 puntos de no-ocurrencia ---igual número que los
eventos inventariados, para mantener el balance de clases--- mediante
muestreo aleatorio con dos criterios de filtrado:
(i)~pendiente superior a 5°, para excluir superficies planas no
susceptibles que aportarían información poco informativa sobre los
condicionantes de la falla, y
(ii)~distancia mínima de 1\,km a cualquier punto del inventario, para
minimizar la autocorrelación espacial entre ocurrencias y
no-ocurrencias.

\subsubsection*{Ajuste del modelo y validación de la dimensión de la base}

El modelo se ajustó con la función \texttt{gam()} del paquete
\texttt{mgcv} \citep{wood2017generalized} en el entorno estadístico R,
especificando distribución binomial, enlace logit y REML como criterio
de selección del suavizado. La dimensión $k$ de cada término de
suavizado determina la complejidad máxima de la curva estimada: valores
de $k$ insuficientes pueden impedir capturar la no-linealidad real de
la relación. La adecuación de $k$ se verificó con la función
\texttt{gam.check()}, que realiza una prueba estadística basada en
residuos por vecinos más próximos: un valor $p$ pequeño ($p < 0{,}05$)
indica que la función aún presenta estructura no capturada con la base
actual y conviene incrementar $k$.

\subsubsection*{Diagnósticos del ajuste: \texttt{gam.check()}}

La función \texttt{gam.check()} produce cuatro gráficos de diagnóstico
que permiten evaluar los supuestos del modelo
(Figura~\ref{fig:gam_check}):

\begin{enumerate}
    \item \textbf{Gráfico Q-Q de residuos de devianza}: compara los
    cuantiles empíricos de los residuos con los cuantiles teóricos de
    la distribución normal. Si los puntos se alinean sobre la recta de
    referencia, los residuos siguen aproximadamente la distribución
    esperada. Desviaciones sistemáticas en los extremos pueden
    indicar valores atípicos o errores de especificación del modelo.

    \item \textbf{Histograma de residuos}: muestra la distribución
    marginal de los residuos de devianza. Se espera una distribución
    aproximadamente simétrica y unimodal centrada en cero. Asimetrías
    marcadas o bimodalidad pueden revelar clases mal representadas o
    datos atípicos.

    \item \textbf{Residuos vs.\ valores ajustados}: verifica la
    homocedasticidad de los residuos. La ausencia de patrones
    sistemáticos (tendencias en forma de embudo, curvatura o
    estructura periódica) confirma que el modelo captura adecuadamente
    la varianza de los datos. Para respuestas binarias es normal
    observar dos franjas de puntos, una por cada clase.

    \item \textbf{Respuesta observada vs.\ ajustada}: contrasta los
    valores medios observados en intervalos del predictor lineal con
    los valores predichos. Una nube de puntos cercana a la diagonal
    confirma que el modelo discrimina correctamente ambas clases en
    todo el rango del predictor.
\end{enumerate}

\subsubsection*{Análisis de concurvidad}

La concurvidad es el análogo no paramétrico de la multicolinealidad:
cuantifica hasta qué punto una función de suavizado $s_j$ puede ser
aproximada por una combinación de las demás funciones del modelo
\citep{wood2017generalized}. Un índice de concurvidad cercano a 1 para
un par de predictores indica que ambos contienen información casi
redundante, lo que compromete la identificabilidad y la precisión de
las curvas estimadas. Los índices de concurvidad se calcularon con la
función \texttt{concurvity()} y se verificó que ningún par de
predictores superara el umbral de $0{,}8$ que se considera indicativo
de concurvidad problemática.

\subsection{Partición y evaluación del desempeño}

El conjunto de datos completo ($n = 10\,974$: 5\,487 ocurrencias y
5\,487 no-ocurrencias) se dividió estratificadamente en 70\,\%
entrenamiento ($n = 7\,682$) y 30\,\% prueba ($n = 3\,292$),
manteniendo la proporción de clases en ambas particiones. El modelo
se ajustó exclusivamente con el subconjunto de entrenamiento y fue
evaluado en el de prueba, el cual no participó en ninguna etapa del
ajuste, garantizando una estimación imparcial del desempeño predictivo.

El umbral óptimo de clasificación (probabilidad predicha por encima de
la cual se asigna la clase ``ocurrencia'') se determinó maximizando el
índice de Youden (TSS) en la curva ROC. Se calcularon las siguientes
métricas sobre el conjunto de prueba:

\begin{itemize}
    \item \textbf{AUC} (\textit{Area Under the ROC Curve}): área bajo
    la curva ROC (Característica Operativa del Receptor). Mide la
    capacidad discriminante del modelo independientemente del umbral
    de clasificación: AUC~$= 0{,}5$ equivale a una predicción aleatoria
    y AUC~$= 1{,}0$ a una discriminación perfecta. Convencionalmente:
    $<0{,}7$ = pobre; $0{,}7$--$0{,}8$ = aceptable;
    $0{,}8$--$0{,}9$ = bueno; $>0{,}9$ = excelente
    \citep{hosmer2013applied}. Es la métrica más empleada en modelos
    de susceptibilidad por ser insensible a la elección del umbral de
    clasificación.

    \item \textbf{TSS} (\textit{True Skill Statistic}, índice de Youden):
    $\mathrm{TSS} = \mathrm{Sensibilidad} + \mathrm{Especificidad} - 1$.
    Equilibra la capacidad del modelo para identificar correctamente
    tanto deslizamientos (sensibilidad) como zonas estables
    (especificidad). Su rango es $[-1,\,1]$: 0 indica rendimiento
    equivalente al azar; valores $> 0{,}6$ se consideran satisfactorios
    para modelos de susceptibilidad. Valores $<0{,}4$ = pobre;
    $0{,}4$--$0{,}6$ = moderado; $0{,}6$--$0{,}8$ = bueno;
    $>0{,}8$ = excelente.

    \item \textbf{Sensibilidad} (recall, tasa de verdaderos positivos):
    fracción de los deslizamientos reales que el modelo clasifica
    correctamente. Una alta sensibilidad reduce el riesgo de no
    identificar zonas peligrosas (falsos negativos), lo que es
    prioritario en aplicaciones de gestión del riesgo.

    \item \textbf{Especificidad} (tasa de verdaderos negativos):
    fracción de las zonas estables que el modelo clasifica
    correctamente. Una alta especificidad evita generar alarmas
    innecesarias sobre terreno seguro (falsos positivos), importante
    para la viabilidad del ordenamiento territorial.

    \item \textbf{F1-Score}: media armónica de precisión
    (fracción de predicciones positivas que son realmente ocurrencias)
    y sensibilidad:
    $F_1 = 2\,\mathrm{TP}/(2\,\mathrm{TP}+\mathrm{FP}+\mathrm{FN})$.
    Útil cuando se busca equilibrar la detección de eventos reales
    con la minimización de falsas alarmas, especialmente en conjuntos
    de datos balanceados.

    \item \textbf{Kappa de Cohen} ($\kappa$): mide la concordancia
    entre clases predichas y observadas corrigiendo por el nivel de
    acuerdo esperado al azar:
    $\kappa = (\mathrm{Po} - \mathrm{Pe})/(1 - \mathrm{Pe})$,
    donde $\mathrm{Po}$ es la precisión observada y $\mathrm{Pe}$
    la esperada por azar. Valores $<0{,}2$ = insignificante;
    $0{,}2$--$0{,}4$ = leve; $0{,}4$--$0{,}6$ = moderado;
    $0{,}6$--$0{,}8$ = sustancial; $>0{,}8$ = casi perfecto
    \citep{cohen1960}.
\end{itemize}

La importancia relativa de cada predictor se cuantificó como la
reducción de devianza al omitir esa variable del modelo completo:
$\Delta D_j = D_{\mathrm{sin}\,j} - D_{\mathrm{completo}}$.
Un valor alto de $\Delta D_j$ indica que el predictor $j$ aporta
información irreemplazable para explicar la distribución espacial
de los movimientos en masa en el Valle de Aburrá.

%===================================================================
\section{Resultados}
%===================================================================

La Figura~\ref{fig:hist} presenta las distribuciones de densidad de
las seis variables para los grupos de deslizamiento y no-ocurrencia.
Las pruebas de Kolmogorov-Smirnov confirmaron diferencias
estadísticamente significativas ($p < 0{,}001$) en todas las variables.

La elevación presenta distribuciones claramente distintas:
los deslizamientos se concentran entre 1\,500 y 2\,000\,m, mientras
las no-ocurrencias son más uniformes. La pendiente muestra
la diferencia más marcada (mayor estadístico $D$ de KS): los
deslizamientos se acumulan en 20°--45° con pico en torno a 30°,
mientras las no-ocurrencias tienen mayor representación en pendientes
suaves. El aspecto evidencia una distribución prácticamente
uniforme para las no-ocurrencias, pero con leve predominio de
deslizamientos en orientaciones oeste-noroeste (225°--315°). Las
curvaturas muestran que los deslizamientos tienden hacia
valores negativos (zonas cóncavas en planta y en perfil), coherente
con la acumulación de agua y material. El TWI muestra la
separación más clara en términos de densidad: los deslizamientos
registran valores consistentemente mayores, confirmando la importancia
de la humedad en la ocurrencia de movimientos en masa en el Valle de Aburrá.

El ajuste del modelo GAM convergió correctamente y los seis términos de suavizado
resultaron altamente significativos ($p < 0{,}001$). La
Figura~\ref{fig:gam_check} presenta los cuatro gráficos de diagnóstico
de \texttt{gam.check()}: el Q-Q de residuos de devianza muestra un
ajuste cercano a la distribución teórica normal; el histograma de
residuos es aproximadamente simétrico; el gráfico de residuos vs.
valores ajustados no evidencia heterocedasticidad sistemática; y el
gráfico respuesta observada vs. ajustada confirma una discriminación
correcta de ambas clases. Las pruebas de adecuación de la dimensión
$k$ no detectaron deficiencias ($p > 0{,}05$ en todos los términos).
El análisis de concurvidad indicó valores moderados ($< 0{,}6$) para
todos los pares de predictores, descartando problemas de
estimabilidad.

\begin{figure}[htbp]
    \centering
    \includegraphics[width=0.85\textwidth]{FIGURAS/gam_check.png}
    \caption{Gráficos de diagnóstico del modelo GAM
    (\texttt{gam.check()}): Q-Q de residuos de devianza
    (sup.\ izq.), histograma de residuos (sup.\ der.),
    residuos vs.\ valores ajustados (inf.\ izq.) y respuesta
    observada vs.\ ajustada (inf.\ der.).}
    \label{fig:gam_check}
\end{figure}

La Figura~\ref{fig:smooths} presenta las seis funciones de suavizado
estimadas, que describen el efecto marginal de cada variable sobre el
predictor lineal $\eta$ (escala logit) manteniendo las demás variables
en sus valores medios.

\begin{figure}[htbp]
    \centering
    \includegraphics[width=1.00\textwidth]{FIGURAS/gam_smooths.png}
    \caption{Funciones de suavizado del GAM para las seis variables
    morfométricas (tres filas, dos columnas). Eje $y$: contribución al
    predictor lineal (logit) centrada en cero. Bandas sombreadas:
    intervalos de confianza al 95\,\%. Puntos grises: residuos de
    devianza parciales. Panel superior derecho: \textit{spline} cúbico
    cíclico del aspecto (\texttt{bs\,=\,"cc"}).}
    \label{fig:smooths}
\end{figure}

\begin{itemize}

\item $s(\text{elevación})$: efecto positivo hasta los 1\,900--2\,000\,m
y posterior descenso. El incremento del logit entre 1\,400 y 1\,900\,m
refleja la concentración de perfiles de meteorización profundos y
urbanización de ladera en esa franja altitudinal; la caída a
elevaciones superiores sugiere zonas con menor densidad
constructiva y/o menor susceptibilidad de los materiales.

\item $s(\text{pendiente})$: presenta un incremento sostenido y pronunciado hasta
los 35°--40°, seguido de una meseta o leve descenso. La forma de la
curva es coherente con la mecánica de la estabilidad de laderas: la
tensión de corte aumenta con $\sin\beta$, pero pendientes muy
pronunciadas tienden a carecer de material removible tras episodios
previos de inestabilidad o erosión.

\item $s_\text{cc}(\text{aspecto})$: Señala un efecto circular con mayor
contribución al logit en orientaciones oeste y noroeste (225°--315°).
El efecto puede estar indicando mayor humedad relativa de las
laderas con esa orientación en el contexto climático del Valle de
Aburrá.

\item $s(\kappa_p\,\text{planiforme})$: efecto negativo en valores
de curvatura planiforme negativos (zonas cóncavas en planta). La
contribución al logit aumenta notablemente para $\kappa_p < -0{,}5$,
confirmando que las vaguadas y quebradas son los sectores de mayor
susceptibilidad por concentración del flujo subsuperficial.

\item $s(\kappa_r\,\text{perfil})$: patrón similar a $\kappa_p$, con
mayor susceptibilidad en zonas cóncavas en perfil ($\kappa_r < 0$),
donde el suelo transportado se acumula y la presión de poros puede alcanzar
valores críticos.

\item $s(\text{TWI})$: presenta un efecto positivo y monótonamente creciente, con
la tasa de incremento más alta entre valores de 3 y 7. Este es el
efecto más pronunciado de todos los predictores, confirmando que la
humedad topográfica es el factor condicionante dominante en el sistema.

\end{itemize}

La Figura~\ref{fig:roc} presenta la curva ROC en el conjunto de prueba.
El AUC de 0{,}908 sitúa el modelo en la categoría de desempeño
\textit{excelente}. La Tabla~\ref{tab:metrics} resume todas las
métricas con el umbral óptimo determinado por el máximo TSS.

\begin{figure}[htbp]
    \centering
    \includegraphics[width=0.65\textwidth]{FIGURAS/curva_roc.png}
    \caption{Curva ROC del modelo GAM evaluada en el conjunto de
    prueba (30\,\%). El punto rojo indica el umbral óptimo (máximo
    TSS/Youden). Línea discontinua: referencia de clasificación
    aleatoria.}
    \label{fig:roc}
\end{figure}

\begin{table}[htbp]
\centering
\caption{Métricas de desempeño del modelo GAM en el conjunto de
prueba con umbral óptimo (máximo TSS).}
\label{tab:metrics}
\begin{tabular}{lcl}
\toprule
Métrica & Valor & Categoría de desempeño \\
\midrule
AUC                        & 0{,}908 & Excelente ($>0{,}9$) \\
TSS (True Skill Statistic) & 0{,}816 & Excelente ($>0{,}8$) \\
Exactitud (Accuracy)       & 0{,}842 & --- \\
Sensibilidad (Recall)      & 0{,}851 & --- \\
Especificidad              & 0{,}833 & --- \\
F1-Score                   & 0{,}847 & --- \\
Kappa de Cohen             & 0{,}684 & Sustancial (0{,}6--0{,}8) \\
\bottomrule
\end{tabular}
\end{table}

El TSS de 0{,}816 indica que el modelo clasifica correctamente el
81{,}6\,\% de los casos más allá de la probabilidad esperada por azar.
La sensibilidad (0{,}851) y la especificidad (0{,}833) son equilibradas,
indicando que el modelo discrimina bien ambas clases sin sesgo hacia
ninguna. El Kappa de 0{,}684 corresponde a acuerdo sustancial.

La Figura~\ref{fig:imp} presenta la reducción de devianza ($\Delta D$)
al omitir cada predictor del modelo completo. El \textbf{TWI} es
claramente el predictor dominante, con una contribución varias veces
superior a la de las demás variables. La \textbf{pendiente} ocupa el
segundo lugar, seguida de \textbf{curvatura de perfil} y
\textbf{curvatura planiforme}. La \textbf{elevación} y el
\textbf{aspecto} tienen contribuciones menores pero estadísticamente
significativas.

\begin{figure}[htbp]
    \centering
    \includegraphics[width=0.72\textwidth]{FIGURAS/importancia_predictores.png}
    \caption{Importancia relativa de los predictores morfométricos
    medida como reducción de devianza ($\Delta D_j$) al excluir cada
    variable del modelo GAM. Mayor $\Delta D_j$ indica mayor
    contribución al poder explicativo.}
    \label{fig:imp}
\end{figure}

La Tabla~\ref{tab:high_susc} presenta los estadísticos de las seis
variables morfométricas para los puntos clasificados con
susceptibilidad alta (probabilidad $\geq 0{,}70$, correspondiente
al cuantil 75 de la distribución de probabilidades predichas).

\begin{table}[htbp]
\centering
\caption{Estadísticos de las variables morfométricas para zonas de
alta susceptibilidad (probabilidad predicha $\geq 0{,}70$). Los valores
caracterizan el rango de condiciones morfométricas asociadas a la
mayor probabilidad de ocurrencia de movimientos en masa.}
\label{tab:high_susc}
\small
\begin{tabular}{lrrrrrr}
\toprule
Variable & Mínimo & Media & Mediana & Máximo & DE & Moda \\
\midrule
Elevación (m)     & 1\,412 & 1\,784 & 1\,745 & 2\,621 & 241 & 1\,698 \\
Pendiente (°)     & 18{,}4 & 33{,}7 & 33{,}1 & 62{,}8 &  8{,}1 & 31{,}4 \\
Aspecto (°)       & 1{,}2  & 242{,}8 & 258{,}3 & 359{,}7 & 98{,}2 & 272{,}1 \\
Curv. planiforme  & $-4{,}12$ & $-0{,}41$ & $-0{,}38$ & 1{,}87 & 0{,}72 & $-0{,}35$ \\
Curv. de perfil   & $-4{,}28$ & $-0{,}39$ & $-0{,}36$ & 1{,}94 & 0{,}76 & $-0{,}31$ \\
TWI               & 4{,}21 & 7{,}14 & 6{,}98 & 12{,}31 & 1{,}64 & 6{,}73 \\
\bottomrule
\end{tabular}
\end{table}

Los valores de la Tabla~\ref{tab:high_susc} caracterizan el perfil
morfométrico de las zonas más susceptibles del Valle de Aburrá:
pendientes pronunciadas (media 33{,}7°), orientación preferentemente
oeste-noroeste (media 242°), curvatura planiforme y de perfil
negativas (convergencia y aceleración), y TWI elevado (media 7{,}14,
claramente superior a la media global de 4{,}87).

La Figura~\ref{fig:mapa_reg} presenta el mapa de susceptibilidad
regional a 50\,m de resolución, con paleta de colores de verde
(baja susceptibilidad) a rojo (alta susceptibilidad). Los patrones
espaciales son geomorfológicamente coherentes: la susceptibilidad
alta (tonos anaranjados y rojos) se concentra en las vertientes
empinadas con vaguadas bien definidas, especialmente en los sectores
de mayor pendiente y convergencia de flujo de ambas vertientes del
valle. El fondo del valle y las divisorias de agua presentan
susceptibilidad baja (tonos verdes). Los resultados son consistentes
con la distribución del inventario mostrada en la
Figura~\ref{fig:inventario}.

\begin{figure}[htbp]
    \centering
    \includegraphics[width=0.72\textwidth]{FIGURAS/mapa_susceptibilidad_regional.png}
    \caption{Mapa de susceptibilidad por movimientos en masa del
    Valle de Aburrá (resolución 50\,m). Paleta verde--amarillo--rojo:
    de baja (0) a alta (1{,}0) susceptibilidad. Se incluyen escala
    gráfica y flecha norte.}
    \label{fig:mapa_reg}
\end{figure}

%===================================================================
\section{Discusión}
%===================================================================

<<<<<<< Updated upstream
Los resultados evidencian que los GAM constituyen una herramienta
metodológicamente robusta para cuantificar relaciones no lineales entre
variables morfométricas e inestabilidad de laderas en el Valle de
Aburrá. En esta sección se discute el desempeño predictivo en
perspectiva comparativa, la interpretación geomorfológica de cada
función de suavizado, las fuentes de incertidumbre del enfoque adoptado
y sus implicaciones para la gestión territorial del riesgo.

%-------------------------------------------------------------------
\subsection*{Desempeño predictivo en perspectiva comparativa}
%-------------------------------------------------------------------
=======
Los resultados confirman que los GAM constituyen una herramienta
metodológicamente idónea para cuantificar el efecto de covariables, en este caso, sobre la ocurrencia de movimientos en masa. El AUC de
0{,}908 y el TSS de 0{,}816 obtenidos en el conjunto de prueba
independiente reflejan la importancia de las variables morfométricas y su inclusión en la evaluación dela susceptibilidad a movimientos en masa.
>>>>>>> Stashed changes

El AUC de 0{,}908 y el TSS de 0{,}816 sitúan el modelo en la
categoría de desempeño excelente y superan el rango típico de los GAM
publicados en contextos de montaña comparables.
\citet{petschko2014} reportaron AUC entre 0{,}72 y 0{,}89 por dominio
litológico en Austria empleando GAM sobre 15\,850\,km²;
\citet{brenning2008statistical} obtuvo AUC $\approx$ 0{,}82 en los
Andes tropicales sudamericanos; y \citet{aristizabal2015landslide}
reportaron valores similares en la región andina colombiana.
El resultado aquí obtenido es, por tanto, superior al rango central
documentado en la literatura para modelos GAM de susceptibilidad.

<<<<<<< Updated upstream
Esta comparación debe interpretarse con cautela, dado que el AUC
es sensible al protocolo de validación, la relación evento/no-evento y
la estrategia de muestreo de ausencias ---factores que difieren entre
estudios \citep{reichenbach2018review}. En particular, la partición
aleatoria estratificada 70/30 utilizada en este trabajo no elimina la
autocorrelación espacial entre los conjuntos de entrenamiento y prueba.
\citet{brenning2012spatial} demostraron que la validación con partición
aleatoria sobreestima el AUC respecto a la validación cruzada espacial
porque los puntos vecinos comparten estructuras locales aprendidas
durante el entrenamiento. \citet{Roberts2017} cuantificaron esta
sobrestimación en 0{,}06--0{,}10 unidades de AUC en conjuntos con
autocorrelación positiva alta. En consecuencia, el AUC = 0{,}908 debe
interpretarse como una medida de ajuste interno al conjunto de datos
disponible, más que como garantía de transferencia espacial a zonas
sin cobertura de inventario. La aplicación de validación cruzada
por bloques espaciales, separando los folds por microcuencas o cuadrantes
disjuntos, permitirá obtener estimaciones más conservadoras y
representativas de la capacidad predictiva real del modelo.

En cuanto a la comparación con algoritmos de \textit{machine learning},
\citet{merghadi2020} concluyeron en una revisión comprensiva de más de
80 estudios que los métodos de ensamble (RF, XGBoost) superan a los
modelos estadísticos clásicos en AUC, alcanzando valores de 0{,}85
a 0{,}95 en conjuntos de prueba aleatorios. Sin embargo, la ventaja
crítica del GAM en el contexto de la gestión del riesgo reside en la
interpretabilidad directa de sus funciones de suavizado: cada
$s_j(\cdot)$ traduce la relación entre variable morfométrica y
probabilidad de falla en un gráfico intuitivo para técnicos y
tomadores de decisión, sin requerir herramientas post-hoc como los
valores SHAP. En entornos regulatorios donde los modelos deben ser
auditables y justificables ante autoridades de planificación territorial,
la interpretabilidad no es una limitación sino un atributo de diseño
deliberado.

%-------------------------------------------------------------------
\subsection*{No-linealidad de las variables morfométricas: evidencia
             y mecanismos}
%-------------------------------------------------------------------

\paragraph{Pendiente y efecto de agotamiento.}
La función $s(\text{pendiente})$ exhibe el rasgo no lineal más
diagnóstico del modelo: un incremento sostenido hasta 35°--40°,
seguido de una meseta y descenso. Este comportamiento, denominado
\textit{efecto de agotamiento} (\textit{exhaustion effect}), ha sido
documentado mediante GAM en contextos tropicales por
\citet{Vorpahl2012}, quienes reportaron una respuesta en forma de
campana con máximo entre 30° y 42° e interpretaron la reducción
posterior como depleción del material susceptible tras ciclos iterativos
de inestabilidad. La revisión de \citet{reichenbach2018review} sobre
565 modelos de susceptibilidad concluye que las respuestas no monótonas
de la pendiente son frecuentes, y que los modelos lineales que las
ignoran sobreestiman sistemáticamente la susceptibilidad en acantilados
rocosos. En el Valle de Aburrá, donde las vertientes más pronunciadas
corresponden a afloramientos de gneis y anfibolitas del Complejo
Cajamarca con escaso perfil de meteorización, la coherencia geomecánica
de este mecanismo es plena: los movimientos en masa en saprolito se
concentran en pendientes moderadas a empinadas (20°--45°), mientras que
las paredes rocosas escarpadas son intrínsecamente estables o han
liberado ya todo su material disponible. Un modelo de regresión logística
que asumiese linealidad en el logit sobreestimaría la susceptibilidad en
acantilados y subestimaría el rango crítico de 30°--35°, produciendo
mapas de zonificación distorsionados respecto a la realidad mecánica del
terreno \citep{sidle2006landslides}.

\paragraph{Elevación: patrón no monotónico e interacción antrópica.}
El pico de susceptibilidad en el intervalo 1\,500--1\,900\,m y el
posterior descenso difieren de las relaciones lineales documentadas en
otros orógenos \citep{guzzetti2006landslide}. Este patrón integra al
menos dos factores concurrentes. En primer lugar, los perfiles de
meteorización más potentes del Valle de Aburrá ---saprolito sobre
esquistos cloríticos del Complejo Cajamarca--- se desarrollan
preferentemente en ese rango altitudinal, donde la combinación de
temperatura, humedad y tiempo de exposición genera horizontes arcillosos
de alta plasticidad sobre los que se producen los deslizamientos
translacionales superficiales. En segundo lugar, la densidad de
asentamientos informales y la presión constructiva sobre ladera alcanzan
su máximo entre 1\,500 y 2\,000\,m, lo que aumenta la probabilidad de
que un evento sea reportado e incorporado al inventario, introduciendo un
sesgo de detección altitudinal que no puede separarse de la señal
geomorfológica real con las variables disponibles. En elevaciones
superiores a 2\,000\,m predominan afloramientos rocosos con menor
espesor de suelo y menor presión antrópica directa, lo que explica la
reducción de susceptibilidad en ese rango.

\paragraph{Aspecto: control mesoclimático y tratamiento circular.}
La mayor susceptibilidad en orientaciones oeste-noroeste (225°--315°)
es consistente con los patrones de circulación atmosférica regional:
los flujos húmedos procedentes del Pacífico generan mayor precipitación
orográfica sobre las vertientes occidentales del Valle de Aburrá.
\citet{aristizabal2016shia} identificaron que el aspecto modula la
humedad antecedente del suelo en cuencas andinas colombianas, condición
previa para la saturación del perfil saprolítico durante tormentas
intensas. En el contexto colombiano, \citet{GomezGutierrez2023}
reportaron que las laderas orientadas al norte y al oeste exhiben mayor
susceptibilidad en la cordillera Oriental, coincidiendo cualitativamente
con los resultados del presente trabajo, aunque los mecanismos climáticos
difieren entre flancos andinos. La literatura señala, sin embargo, que
el efecto del aspecto es altamente dependiente del contexto climático
local y no admite generalización \citep{reichenbach2018review}.

Desde el punto de vista metodológico, el uso del \textit{spline} cúbico
cíclico para el aspecto representa una mejora respecto a los enfoques
habituales que descomponen el aspecto en componentes seno/coseno
---los cuales suponen implícitamente una respuesta sinusoidal de período
fijo--- o que lo categorizan en clases discretas
\citep{wilson2000terrain}. La continuidad garantizada en 0°/360° evita
el artefacto numérico que los TPS estándar introducen al tratar ese
punto como discontinuidad, y la curva periódica resultante captura
asimetrías entre vertientes de manera flexible con base estadística.

\paragraph{Dominancia del TWI y advertencia metodológica.}
El TWI emerge como el predictor dominante por amplio margen,
confirmando la primacía del control hidrológico sobre la inestabilidad
superficial en este ambiente. Este resultado converge con lo reportado
por \citet{Munoz2023heliyon} en los Andes ecuatorianos, donde el TWI
y la curvatura planiforme explicaron conjuntamente la mayor parte de la
varianza de ocurrencia, y con \citet{moralesvelez2024}, quienes
demostraron en la región andina colombiana que la humedad del suelo
---expresada a través de índices topográficos equivalentes al TWI---
media la relación entre el forzamiento del ENOS y la frecuencia de
deslizamientos. Desde una perspectiva física, la dominancia del TWI
refleja que la acumulación de presión de poros en zonas de convergencia
topográfica es el mecanismo desencadenante dominante de los movimientos
en masa superficiales en suelos saprolíticos bajo precipitación intensa
\citep{beven1979physically}.

Sin embargo, es necesario señalar una advertencia metodológica
importante. En este trabajo se empleó la aproximación local
$\mathrm{TWI} \approx \ln(1/\tan\beta)$, que es una transformación
logarítmica monotónica de la pendiente. Esta simplificación, adoptada
por limitaciones de cómputo del área de contribución $A_s$ a alta
resolución, hace que la pendiente y el TWI sean funcionalmente
colineales en el modelo, aunque la concurvidad medida resultó moderada
($< 0{,}6$). La dominancia reportada del TWI en el análisis de
reducción de devianza podría estar amplificada por este efecto: parte
de la señal atribuida al TWI correspondería en realidad a la componente
logarítmica de la pendiente que el modelo GAM no puede separar con la
aproximación empleada. El cálculo del TWI con acumulación D-infinito
\citep{tarboton1997}, que incorpora el área de contribución aguas arriba
completa, es la principal mejora metodológica requerida para versiones
futuras del modelo.

\paragraph{Curvaturas planiforme y de perfil.}
La mayor susceptibilidad en zonas cóncavas (curvaturas negativas)
responde al mecanismo físico de concentración del flujo subsuperficial
en vaguadas: la convergencia lateral aumenta la presión de poros y
favorece la licuación del saprolito durante eventos extremos
\citep{beven1979physically}. Las funciones $s(\kappa_p)$ y $s(\kappa_r)$
exhiben asimetría marcada: el efecto negativo en zonas cóncavas es más
pronunciado que el efecto positivo en convexas, lo que sugiere que la
concentración de flujo actúa como umbral de desestabilización en lugar
de como continuum lineal. Este resultado es consistente con la
distribución del inventario (Figura~\ref{fig:inventario}), que muestra
concentración de eventos en las quebradas y vaguadas del sistema hídrico
del Valle de Aburrá.

%-------------------------------------------------------------------
\subsection*{Sesgo del inventario y propagación al modelo}
%-------------------------------------------------------------------

El inventario del SIATA es el más completo disponible para el Valle de
Aburrá, pero comparte los sesgos estructurales documentados en
inventarios basados en reportes de emergencia: sobrerepresentación de
eventos en zonas urbanizadas y accesibles, y subregistro en vertientes
forestadas y cuencas remotas. \citet{steger2017} demostraron que
incluso inventarios con este patrón producen modelos con AUC
aparentemente alto ($> 0{,}80$), pero con relaciones predictor-respuesta
distorsionadas que favorecen variables correlacionadas con la presencia
humana. \citet{Steger2021correlation} confirmaron que la correlación
estadística entre predictores morfométricos y deslizamientos puede
reflejar el patrón de recolección de datos más que la susceptibilidad
intrínseca del terreno, fenómeno que denominaron sesgo de accesibilidad.
En el Valle de Aburrá, el pico de susceptibilidad en el rango altitudinal
1\,500--1\,900\,m podría estar parcialmente condicionado por la mayor
probabilidad de registro en ese rango, donde se concentra la
infraestructura de monitoreo y respuesta del SIATA. Diferenciar la
señal geomorfológica real del artefacto de inventario requeriría un
análisis de sensibilidad que compare versiones del modelo entrenadas
sobre subconjuntos del inventario estratificados por accesibilidad y
distancia a infraestructura, una extensión necesaria antes de utilizar
el mapa para zonificación definitiva.

%-------------------------------------------------------------------
\subsection*{Parsimonia morfométrica y variables omitidas}
%-------------------------------------------------------------------

La decisión de restringir el modelo a variables morfométricas derivadas
del DEM responde a un criterio de consistencia de escala: el DEM
aerofotogramétrico de 2\,m es la fuente de mayor precisión y
homogeneidad disponible, mientras que la geología regional (1:50\,000)
y el uso del suelo presentan incertidumbres espaciales de decenas de
metros incompatibles con el píxel fino del modelo. \citet{reichenbach2018review}
documentaron que, en el 40\,\% de los estudios revisados, la morfometría
topográfica explica más del 60\,\% de la varianza del modelo de
susceptibilidad, validando el enfoque morfométrico como punto de
partida autosuficiente.

No obstante, la omisión de la litología implica que los efectos
estimados de la elevación y las curvaturas absorben parte de la
variabilidad litológica no capturada explícitamente: el pico de
susceptibilidad en 1\,500--1\,900\,m puede estar capturando
simultáneamente la franja donde afloran los esquistos más meteorizados
y la zona de mayor densidad constructiva. La omisión de la precipitación
---detonante principal de los eventos del inventario
\citep{moralesvelez2024}--- limita la interpretación causal del
modelo: el TWI actúa como proxy estático de la propensión
topográfica a la saturación pero no captura la dinámica temporal de
recarga. Esta restricción impide la transición del mapa de
susceptibilidad estático a un sistema de alerta temprana con resolución
temporal, extensión de alto valor operativo para el SIATA.

La resolución del DEM tiene un efecto demostrado sobre las variables
morfométricas derivadas. \citet{Hong2019scale} mostraron que las
variables de curvatura exhiben la mayor sensibilidad a la resolución,
con variaciones pronunciadas del factor de certeza entre DEMs de 5 y
30\,m, mientras que el TWI muestra sensibilidad moderada. En este
trabajo se derivaron las variables morfométricas a 2\,m y el TWI a
10\,m; la coherencia entre ambas escalas de análisis y su efecto sobre
los resultados del GAM merece exploración sistemática, especialmente
para evaluar si el ruido micro-topográfico del DEM de 2\,m introduce
variabilidad artefactual en las curvaturas que distorsione las funciones
de suavizado estimadas.
=======
La pendiente ejerce el segundo efecto más importante (Figura~\ref{fig:imp}).
La función $s(\text{slope})$ muestra un incremento casi monotónico
hasta los 35°--40°, coherente con el aumento de la componente
tangencial del peso que genera la tensión de corte en el plano de
falla. La reducción de la susceptibilidad predicha para pendientes
muy elevadas ($>45^{\circ}$) es un fenómeno conocido como \textit{efecto de
agotamiento}: en laderas extrem\-adamente pronunciadas, los episodios
sucesivos de inestabilidad han removido el material susceptible,
dejando afloramientos rocosos competentes.

El \textit{spline} cíclico permite identificar que las laderas con
orientación oeste-noroeste (225°--315°) presentan mayor susceptibilidad.
En el Valle de Aburrá, estas orientaciones corresponden a las vertientes...

Las curvaturas planiforme y de perfil aportan información
complementaria sobre la geometría tridimensional de la ladera. La
curvatura planiforme negativa (zonas cóncavas en planta, vaguadas)
indica convergencia del flujo subsuperficial, generando presiones de
poros que pueden reducir la resistencia al corte del material hasta
valores críticos. La curvatura de perfil negativa (zonas cóncavas en
perfil) indica acumulación de material residual y aceleración del
flujo, incrementando tanto el peso del material como la energía
cinética disponible para la movilización. La inclusión conjunta de
ambas curvaturas mejora significativamente la representación de la
inestabilidad local respecto a modelos con solo pendiente y elevación.

El TWI es el predictor dominante (Figura~\ref{fig:imp}), con un efecto
positivo pronunciado sobre la susceptibilidad. Esto es
consistente con el mecanismo de falla predominante en el Valle de
Aburrá: deslizamientos superficiales detonados
por la saturación de los horizontes de meteorización durante episodios
de lluvia intensa d earga duración. La tabla~\ref{tab:high_susc} confirma que las zonas
de alta susceptibilidad presentan valores de TWI medios de 7{,}14,
frente a la media global de 4{,}87, una diferencia de 47\,\%.
>>>>>>> Stashed changes

%===================================================================
\section{Conclusiones}
%===================================================================

\begin{enumerate}
    \item Los GAM con regresión logística y \textit{thin plate splines}
          permiten cuantificar el efecto no lineal e interpretable de
          variables morfométricas sobre la ocurrencia de movimientos en
          masa, con AUC = 0{,}908 y TSS = 0{,}816 en el Valle de
          Aburrá.

    \item El TWI es el predictor dominante, seguido de la pendiente
          y las curvaturas topográficas, lo que resalta el papel de la
          humedad topográfica y la geometría tridimensional de la ladera
          como factores condicionantes principales.

    \item Las funciones de suavizado identifican rangos críticos de
          cada variable: pendientes 20°--45°, elevaciones 1\,500--2\,000\,m,
          orientaciones oeste-noroeste, zonas cóncavas en planta y en
          perfil, y TWI $> 5$.

    \item El tratamiento circular del aspecto mediante \textit{spline}
          cúbico cíclico, la inclusión de curvaturas y la evaluación con
          umbral optimizado por TSS en datos de prueba independientes son
          mejoras metodológicas que deben adoptarse sistemáticamente en
          modelación de susceptibilidad.

    \item El mapa de susceptibilidad a 50\,m y las tablas de rangos
          críticos de variables morfométricas son insumos directamente
          utilizables para la zonificación del riesgo y el ordenamiento
          territorial en el Área Metropolitana del Valle de Aburrá.
\end{enumerate}

%===================================================================
\section*{Referencias}
%===================================================================

\bibliographystyle{plainnat}
\bibliography{referencias}

\end{document}
